\chapter{Chapter09 — Friedgut Necessity}
%%% generated by the blueprint pipeline from FABL.Chapter09.FriedgutNecessity
\section{Overview}

Declarations of the Lean module \texttt{FABL.Chapter09.FriedgutNecessity}.

\section{Declarations}

\begin{definition}[Sign cube address bijection]
\lean{FABL.addressingAddressEquiv}
\label{FABL.addressingAddressEquiv}
\leanok
\uses{FABL.Sign, FABL.SignCube}
For each natural number $m$, there is a bijection between the set of length-$m$ vectors
with entries in $\{-1,1\}$ and the set of integers $0,1,\ldots,2^m-1$.
\end{definition}

\begin{definition}[Addressing input decomposition]
\lean{FABL.addressingInputEquiv}
\label{FABL.addressingInputEquiv}
\leanok
\uses{FABL.Sign, FABL.SignCube}
For each natural number $m$, there is a bijection between sign vectors of length $m+2^m$
and pairs consisting of a sign vector of length $m$ and a sign vector indexed by the
integers $0,1,\ldots,2^m-1$.
\end{definition}

\begin{definition}[Addressing function]
\lean{FABL.addressingFunction}
\label{FABL.addressingFunction}
\leanok
\uses{FABL.BooleanFunction, FABL.addressingAddressEquiv, FABL.addressingInputEquiv}
For each natural number $m$, define a function on sign vectors of length $m+2^m$ as
follows: interpret the first $m$ coordinates as an integer in $\{0,\ldots,2^m-1\}$, then
output the data coordinate indexed by that integer.
\end{definition}

\begin{definition}[Address coordinate embedding]
\lean{FABL.addressingAddressCoordinate}
\label{FABL.addressingAddressCoordinate}
\leanok
For each natural number $m$ and each index $i\in\{0,\ldots,m-1\}$, the corresponding
address coordinate is the coordinate with index $i$ among the first $m$ positions of a
vector of length $m+2^m$.
\end{definition}

\begin{definition}[Data coordinate embedding]
\lean{FABL.addressingDataCoordinate}
\label{FABL.addressingDataCoordinate}
\leanok
For each natural number $m$ and each index $j\in\{0,\ldots,2^m-1\}$, the corresponding
data coordinate is the coordinate with index $j$ among the final $2^m$ positions of a
vector of length $m+2^m$.
\end{definition}

\begin{theorem}[Address coordinates under decomposition]
\lean{FABL.addressingInputEquiv_addressCoordinate}
\label{FABL.addressingInputEquiv_addressCoordinate}
\leanok
\uses{FABL.Sign, FABL.SignCube, FABL.addressingAddressCoordinate, FABL.addressingInputEquiv}
\difficulty{3}
For every natural number $m$, every sign vector $x$ of length $m+2^m$, and every address
index $i<m$, the $i$th coordinate of the address part in the decomposition of $x$ equals
the coordinate of $x$ at address position $i$.
\end{theorem}

\begin{theorem}[Data coordinates under decomposition]
\lean{FABL.addressingInputEquiv_dataCoordinate}
\label{FABL.addressingInputEquiv_dataCoordinate}
\leanok
\uses{FABL.Sign, FABL.SignCube, FABL.addressingDataCoordinate, FABL.addressingInputEquiv}
\difficulty{3}
For every natural number $m$, every sign vector $x$ of length $m+2^m$, and every data
index $j<2^m$, the $j$th coordinate of the data part in the decomposition of $x$ equals
the coordinate of $x$ at data position $j$.
\end{theorem}

\begin{definition}[Selected data coordinate]
\lean{FABL.addressingSelectedCoordinate}
\label{FABL.addressingSelectedCoordinate}
\leanok
\uses{FABL.SignCube, FABL.addressingAddressEquiv, FABL.addressingDataCoordinate, FABL.addressingInputEquiv}
For each natural number $m$ and sign vector $x$ of length $m+2^m$, interpret its first
$m$ coordinates as an integer in $\{0,\ldots,2^m-1\}$; the selected coordinate is the
corresponding data coordinate.
\end{definition}

\begin{theorem}[Address preservation under a selected flip]
\lean{FABL.addressingInputEquiv_first_flip_selected}
\label{FABL.addressingInputEquiv_first_flip_selected}
\leanok
\uses{FABL.SignCube, FABL.addressingAddressCoordinate, FABL.addressingInputEquiv, FABL.addressingInputEquiv_addressCoordinate, FABL.addressingSelectedCoordinate, FABL.flipCoordinate, FABL.setCoordinate_apply_of_ne}
\difficulty{4}
For every natural number $m$ and every sign vector $x$ of length $m+2^m$, flipping the
coordinate indexed by the address encoded by the first component of the decomposition of
$x$ leaves that first component unchanged.
\end{theorem}

\begin{theorem}[Selected coordinate is flip-invariant]
\lean{FABL.addressingSelectedCoordinate_flip}
\label{FABL.addressingSelectedCoordinate_flip}
\leanok
\uses{FABL.SignCube, FABL.addressingAddressEquiv, FABL.addressingDataCoordinate, FABL.addressingInputEquiv_first_flip_selected, FABL.addressingSelectedCoordinate, FABL.flipCoordinate}
\difficulty{3}
For every natural number $m$ and every sign vector $x$ of length $m+2^m$, the data
coordinate selected by the address coordinates remains the same after that selected
coordinate is flipped.
\end{theorem}

\begin{theorem}[Flipping the selected data coordinate reverses the output]
\lean{FABL.addressingFunction_flip_selected}
\label{FABL.addressingFunction_flip_selected}
\leanok
\uses{FABL.SignCube, FABL.addressingAddressEquiv, FABL.addressingDataCoordinate, FABL.addressingFunction, FABL.addressingInputEquiv, FABL.addressingInputEquiv_dataCoordinate, FABL.addressingInputEquiv_first_flip_selected, FABL.addressingSelectedCoordinate, FABL.flipCoordinate, FABL.setCoordinate}
\difficulty{4}
Let $m$ be a natural number and let $x$ be a sign vector with $m+2^m$ coordinates,
divided into $m$ address coordinates and $2^m$ data coordinates. Use a bijection from
sign vectors of length $m$ to the indices $\in{2^m}$ to select an index from the address
coordinates; the addressing function returns the data sign at that index. Flipping the
data coordinate selected by $x$ changes the addressing function's value at $x$ to its
negative.
\end{theorem}

\begin{theorem}[Data flip preserves address part]
\lean{FABL.addressingInputEquiv_first_flip_data}
\label{FABL.addressingInputEquiv_first_flip_data}
\leanok
\uses{FABL.SignCube, FABL.addressingAddressCoordinate, FABL.addressingDataCoordinate, FABL.addressingInputEquiv, FABL.addressingInputEquiv_addressCoordinate, FABL.flipCoordinate, FABL.setCoordinate_apply_of_ne}
\difficulty{4}
For every natural number $m$, every sign vector $x$ of length $m+2^m$, and every data
index $j<2^m$, flipping coordinate $j$ leaves the address part of the decomposition of
$x$ unchanged.
\end{theorem}

\begin{theorem}[Irrelevant data-coordinate flip]
\lean{FABL.addressingFunction_flip_data_of_ne}
\label{FABL.addressingFunction_flip_data_of_ne}
\leanok
\uses{FABL.SignCube, FABL.addressingAddressEquiv, FABL.addressingDataCoordinate, FABL.addressingFunction, FABL.addressingInputEquiv, FABL.addressingInputEquiv_dataCoordinate, FABL.addressingInputEquiv_first_flip_data, FABL.flipCoordinate, FABL.setCoordinate_apply_of_ne}
\difficulty{4}
For every natural number $m$, every sign vector $x$ of length $m+2^m$, and every data
index $j<2^m$ different from the index encoded by the address coordinates of $x$,
flipping the data coordinate indexed by $j$ leaves the addressing function's value at
$x$ unchanged.
\end{theorem}

\begin{definition}[Relevant addressing coordinates]
\lean{FABL.addressingRelevantCoordinates}
\label{FABL.addressingRelevantCoordinates}
\leanok
\uses{FABL.SignCube, FABL.addressingAddressCoordinate, FABL.addressingSelectedCoordinate}
For each natural number $m$ and sign vector $x$ of length $m+2^m$, the relevant
coordinates are all of the first $m$ address coordinates together with the single data
coordinate selected by the address in $x$.
\end{definition}

\begin{theorem}[Pivotal coordinates of the addressing function]
\lean{FABL.isPivotal_addressingFunction_mem_relevant}
\label{FABL.isPivotal_addressingFunction_mem_relevant}
\leanok
\uses{FABL.IsPivotal, FABL.SignCube, FABL.addressingAddressCoordinate, FABL.addressingAddressEquiv, FABL.addressingDataCoordinate, FABL.addressingFunction, FABL.addressingFunction_flip_data_of_ne, FABL.addressingInputEquiv, FABL.addressingRelevantCoordinates, FABL.addressingSelectedCoordinate}
\difficulty{5}
Let $m\in\mathbb N$, and consider the addressing function on sign vectors with $m$
address coordinates and $2^m$ data coordinates: the address selects one data coordinate,
whose sign is the output. For any input $x$ and coordinate $i$, if changing coordinate
$i$ changes the output, then $i$ is either an address coordinate or the data coordinate
selected by the address of $x$.
\end{theorem}

\begin{theorem}[Bound on Relevant Coordinates]
\lean{FABL.card_addressingRelevantCoordinates_le}
\label{FABL.card_addressingRelevantCoordinates_le}
\leanok
\uses{FABL.SignCube, FABL.addressingAddressCoordinate, FABL.addressingRelevantCoordinates, FABL.addressingSelectedCoordinate}
\difficulty{4}
Let $m$ be a natural number and let $x$ be a sign assignment on $m+2^m$ coordinates. The
set consisting of the $m$ address coordinates and the data coordinate selected by the
address in $x$ has cardinality at most $m+1$.
\end{theorem}

\begin{theorem}[Sensitivity of the addressing function]
\lean{FABL.sensitivity_addressingFunction_le}
\label{FABL.sensitivity_addressingFunction_le}
\leanok
\uses{FABL.SignCube, FABL.addressingFunction, FABL.card_addressingRelevantCoordinates_le, FABL.isPivotal_addressingFunction_mem_relevant, FABL.sensitivity}
\difficulty{3}
For every natural number $m$ and every sign vector consisting of an $m$-coordinate
address followed by a table of $2^m$ signs, the addressing function's sensitivity at
that vector is at most $m+1$. Here sensitivity counts the coordinates whose individual
sign flip changes the function's output; the addressing function outputs the table entry
selected by the address.
\end{theorem}

\begin{theorem}[Addressing function influence bound]
\lean{FABL.totalInfluence_addressingFunction_le}
\label{FABL.totalInfluence_addressingFunction_le}
\leanok
\uses{FABL.BooleanFunction.toReal, FABL.SignCube, FABL.addressingFunction, FABL.sensitivity, FABL.sensitivity_addressingFunction_le, FABL.totalInfluence, FABL.totalInfluence_toReal_eq_expect_sensitivity}
\difficulty{3}
For every natural number $m$, the total influence of the addressing function on sign
vectors of dimension $m+2^m$ is at most $m+1$.
\end{theorem}

\begin{definition}[Selected coordinate flip bijection]
\lean{FABL.addressingFlipEquiv}
\label{FABL.addressingFlipEquiv}
\leanok
\uses{FABL.SignCube, FABL.addressingSelectedCoordinate, FABL.addressingSelectedCoordinate_flip, FABL.flipCoordinate, FABL.flipCoordinate_flipCoordinate}
For each natural number $m$, define a bijection on sign vectors of length $m+2^m$ by
flipping the sign at the data coordinate selected by the vector's address.
\end{definition}

\begin{theorem}[Addressing flip map]
\lean{FABL.addressingFlipEquiv_apply}
\label{FABL.addressingFlipEquiv_apply}
\leanok
\uses{FABL.SignCube, FABL.addressingFlipEquiv, FABL.addressingSelectedCoordinate, FABL.flipCoordinate}
\difficulty{1}
For every natural number $m$ and every sign vector $x$ of length $m+2^m$, the addressing
flip map sends $x$ to the vector obtained by flipping the data coordinate selected by
the address encoded in the first $m$ coordinates of $x$.
\end{theorem}

\begin{theorem}[Addressing flip outside the dependence set]
\lean{FABL.eq_addressingFlip_of_dependsOn}
\label{FABL.eq_addressingFlip_of_dependsOn}
\leanok
\uses{FABL.BooleanFunction, FABL.SignCube, FABL.addressingFlipEquiv, FABL.addressingFlipEquiv_apply, FABL.addressingSelectedCoordinate, FABL.flipCoordinate, FABL.setCoordinate_apply_of_ne}
\difficulty{3}
Let $m$ be a natural number, let $J$ be a set of coordinates in a sign cube of dimension
$m+2^m$, and let $g$ be a Boolean function on that cube that depends only on coordinates
in $J$. For any input $x$ whose address-selected coordinate is not in $J$, flipping that
selected coordinate leaves the value of $g$ unchanged.
\end{theorem}

\begin{definition}[Addressing error indicator]
\lean{FABL.addressingErrorIndicator}
\label{FABL.addressingErrorIndicator}
\leanok
\uses{FABL.BooleanFunction, FABL.SignCube}
For Boolean functions $f,g\colon\{−1,1\}^N\to\{−1,1\}$ and an input $x\in\{−1,1\}^N$,
define their addressing error indicator to be $1$ if $f(x)\ne g(x)$ and $0$ otherwise.
\end{definition}

\begin{theorem}[Addressing error under a selected-coordinate flip]
\lean{FABL.one_le_addressingErrorIndicator_add_flip}
\label{FABL.one_le_addressingErrorIndicator_add_flip}
\leanok
\uses{FABL.BooleanFunction, FABL.SignCube, FABL.addressingErrorIndicator, FABL.addressingFlipEquiv, FABL.addressingFlipEquiv_apply, FABL.addressingFunction, FABL.addressingFunction_flip_selected, FABL.addressingSelectedCoordinate, FABL.eq_addressingFlip_of_dependsOn}
\difficulty{5}
Let $m$ be a nonnegative integer, let $J$ be a set of coordinates in the cube
$\{−1,1\}^{m+2^m}$, and let $g:\{−1,1\}^{m+2^m}\to\{−1,1\}$ depend only on the
coordinates in $J$. For any input $x$ whose address-selected coordinate lies outside
$J$, the sum of the error indicators of the addressing function and $g$, evaluated at
$x$ and at the input obtained by flipping that selected coordinate, is at least $1$.
\end{theorem}

\begin{theorem}[Error count under the addressing flip]
\lean{FABL.sum_addressingErrorIndicator_flip}
\label{FABL.sum_addressingErrorIndicator_flip}
\leanok
\uses{FABL.BooleanFunction, FABL.addressingErrorIndicator, FABL.addressingFlipEquiv, FABL.addressingFunction}
\difficulty{1}
Let $m$ be a natural number and let $g$ be a Boolean function on the sign cube with
$m+2^m$ coordinates. The sum, over all inputs, of the error indicator comparing the
addressing function with $g$ after applying the addressing flip equals the corresponding
error-indicator sum before the flip.
\end{theorem}

\begin{theorem}[Addressing error and unselected coordinates]
\lean{FABL.two_mul_sum_addressingErrorIndicator_ge_event_card}
\label{FABL.two_mul_sum_addressingErrorIndicator_ge_event_card}
\leanok
\uses{FABL.BooleanFunction, FABL.SignCube, FABL.addressingErrorIndicator, FABL.addressingFlipEquiv, FABL.addressingFunction, FABL.addressingSelectedCoordinate, FABL.one_le_addressingErrorIndicator_add_flip, FABL.sum_addressingErrorIndicator_flip}
\difficulty{5}
Let $m$ be a natural number, let $J$ be a set of coordinates in a cube of dimension
$m+2^m$, and let $g$ be a sign-valued Boolean function on that cube which depends only
on coordinates in $J$. The sum, over all inputs $x$, of the indicator that $g(x)$
differs from the addressing function is at least half the number of inputs whose
selected coordinate is not in $J$.
\end{theorem}

\begin{definition}[Addressing data indices selected by a set]
\lean{FABL.addressingDataIndicesIn}
\label{FABL.addressingDataIndicesIn}
\leanok
\uses{FABL.addressingDataCoordinate}
For a natural number $m$ and a set $J$ of coordinates in the finite index set of size
$m+2^m$, define the selected data indices to be those $j\in\{0,\ldots,2^m-1\}$ whose
addressing data coordinate belongs to $J$.
\end{definition}

\begin{theorem}[Data Indices in a Coordinate Set]
\lean{FABL.card_addressingDataIndicesIn_le}
\label{FABL.card_addressingDataIndicesIn_le}
\leanok
\uses{FABL.addressingDataCoordinate, FABL.addressingDataIndicesIn}
\difficulty{4}
Let $m$ be a natural number and let $J$ be a subset of the $m+2^m$ coordinates. The
number of indices $j<2^m$ whose corresponding data coordinate belongs to $J$ is at most
$|J|$.
\end{theorem}

\begin{theorem}[Data Indices Outside a Coordinate Set]
\lean{FABL.card_addressingDataIndicesOutside_ge}
\label{FABL.card_addressingDataIndicesOutside_ge}
\leanok
\uses{FABL.addressingDataCoordinate, FABL.card_addressingDataIndicesIn_le}
\difficulty{4}
Let $m>0$ and let $J$ be a subset of the $m+2^m$ coordinates with $|J|\le 2^{m-1}$. At
least $2^{m-1}$ indices $j<2^m$ have their corresponding data coordinate outside $J$.
\end{theorem}

\begin{theorem}[Probability of selecting outside a coordinate set]
\lean{FABL.uniformProbability_addressingSelectedOutside}
\label{FABL.uniformProbability_addressingSelectedOutside}
\leanok
\uses{FABL.Sign, FABL.SignCube, FABL.addressingAddressEquiv, FABL.addressingDataCoordinate, FABL.addressingInputEquiv, FABL.addressingSelectedCoordinate, FABL.uniformProbability}
\difficulty{5}
Let $m$ be a natural number and let $J$ be a set of coordinates among the first $m+2^m$
positions. Choose uniformly a sign vector on these positions, use its first $m$ signs to
select one of the $2^m$ data positions, and let the selected coordinate be the
corresponding data coordinate. The probability that this coordinate lies outside $J$
equals the number of data positions outside $J$, divided by $2^m$.
\end{theorem}

\begin{theorem}[Lower bound for selecting outside a coordinate set]
\lean{FABL.uniformProbability_addressingSelectedOutside_ge_half}
\label{FABL.uniformProbability_addressingSelectedOutside_ge_half}
\leanok
\uses{FABL.SignCube, FABL.addressingDataCoordinate, FABL.addressingSelectedCoordinate, FABL.card_addressingDataIndicesOutside_ge, FABL.uniformProbability, FABL.uniformProbability_addressingSelectedOutside}
\difficulty{4}
Let $m$ be a positive natural number and let $J$ be a set of coordinates among the first
$m+2^m$ positions, with $|J|\le 2^{m-1}$. Choose uniformly a sign vector on these
positions and use its first $m$ signs to select one of the $2^m$ data positions. The
probability that the selected data coordinate lies outside $J$ is at least $1/2$.
\end{theorem}

\begin{theorem}[Addressing function error bound]
\lean{FABL.two_mul_relativeHammingDist_addressing_ge_event}
\label{FABL.two_mul_relativeHammingDist_addressing_ge_event}
\leanok
\uses{FABL.BooleanFunction, FABL.SignCube, FABL.addressingErrorIndicator, FABL.addressingFunction, FABL.addressingSelectedCoordinate, FABL.relativeHammingDist, FABL.two_mul_sum_addressingErrorIndicator_ge_event_card, FABL.uniformProbability, FABL.uniformProbability_ne_eq_relativeHammingDist}
\difficulty{4}
Let $m$ be a nonnegative integer. Encode an input as an address in $\{-1,1\}^m$ together
with one sign-valued data bit for each of the $2^m$ addresses, and let $A_m$ return the
data bit at the encoded address. For a set $J$ of coordinates among these $m+2^m$ input
coordinates, let $g:\{-1,1\}^{m+2^m}\to\{-1,1\}$ depend only on coordinates in $J$.
Under the uniform input distribution, the probability that the coordinate selected by
the address is outside $J$ is at most twice the fraction of inputs on which $A_m$ and
$g$ disagree.
\end{theorem}

\begin{theorem}[Addressing function distance lower bound]
\lean{FABL.relativeHammingDist_addressing_ge_quarter}
\label{FABL.relativeHammingDist_addressing_ge_quarter}
\leanok
\uses{FABL.BooleanFunction, FABL.addressingFunction, FABL.relativeHammingDist, FABL.two_mul_relativeHammingDist_addressing_ge_event, FABL.uniformProbability_addressingSelectedOutside_ge_half}
\difficulty{3}
For every positive integer $m$, let $J$ be a set of coordinates among the $m+2^m$ input
coordinates with $|J|\le 2^{m-1}$. If a Boolean function $g$ on the sign cube of
dimension $m+2^m$ depends only on coordinates in $J$, then its relative Hamming distance
from the addressing function is at least $1/4$.
\end{theorem}

\begin{theorem}[Distance from the addressing function to juntas]
\lean{FABL.not_close_addressing_of_isKJunta}
\label{FABL.not_close_addressing_of_isKJunta}
\leanok
\uses{FABL.BooleanFunction, FABL.IsKJunta, FABL.addressingFunction, FABL.relativeHammingDist, FABL.relativeHammingDist_addressing_ge_quarter}
\difficulty{4}
Let $m$ be a positive integer, and let $A_m$ be the Boolean function on $m+2^m$ signs
that uses the first $m$ signs to select one of the remaining $2^m$ signs, using any
fixed bijection between the $m$-sign strings and the $2^m$ positions. If a Boolean
function $g$ on the same inputs depends on at most $2^{m-1}$ coordinates, then its
relative Hamming distance from $A_m$ is at least $1/4$.
\end{theorem}

\begin{theorem}[Addressing function witness bounds]
\lean{FABL.exercise9_26_addressing_witness}
\label{FABL.exercise9_26_addressing_witness}
\leanok
\uses{FABL.BooleanFunction, FABL.BooleanFunction.toReal, FABL.IsKJunta, FABL.addressingFunction, FABL.not_close_addressing_of_isKJunta, FABL.relativeHammingDist, FABL.totalInfluence, FABL.totalInfluence_addressingFunction_le}
\difficulty{1}
For every positive integer $m$, the addressing function on $m+2^m$ sign coordinates has
total influence at most $m+1$. Moreover, every Boolean function on those coordinates
that depends on at most $2^{m-1}$ coordinates has relative Hamming distance at least
$1/4$ from the addressing function.
\end{theorem}
